\documentclass{article}
\usepackage[T1]{fontenc}
\usepackage{lmodern}
\usepackage{graphicx}
\usepackage{amsmath,amssymb}
\usepackage[a4paper,margin=2cm]{geometry}
\usepackage[absolute,overlay]{textpos}
\setlength{\TPHorizModule}{1mm}
\setlength{\TPVertModule}{1mm}
\usepackage[indonesian]{babel}
\usepackage{hyperref}
\setlength{\emergencystretch}{2em}
% unit: Halaman 32

\begin{document}

% === Halaman 32 (llm) ===

\textcolor{red}{Plainteks}: \begin{center} \begin{tabular}{|p{0.7\textwidth}|} \hline Setelah Presiden Venezuela Nicolas Maduro digulingkan dan ditahan di New York atas tuduhan terorisme narkoba, kepemimpinan di Caracas kini berpindah tangan ke pemerintahan sementara yang dikawal Amerika Serikat (AS). Meski demikian, jaringan loyalis dan aparat keamanan yang menopang kekuasaan Maduro tetap aktif dan memiliki pengaruh besar. Bahkan, tercatat kini ada sepuluh orang yang memegang kendali negara kaya minyak itu, baik dari kubu Venezuela maupun AS. \\ \hline \end{tabular} \end{center} \vspace{10mm} \textcolor{red}{Cipherteks}: \begin{center} \begin{tabular}{|p{0.7\textwidth}|} \hline Ti2TVWlery8mUA6rmIz0IC17mXmlz1ykQH2D+rc5fgQyJkFyNfKcNbnOeJhS+s Ag3wyyPuiuR4xjdpvj0h7jQnzdjcLOj++ /jm6enh7KO6zOvg2RgVdgNcFyaUZZiZ/x+v /2ZmeYUrACii7TgdDb84ypbvSSwSZ32SuOn0JFurxB9ZUpSRHx8b17ff8bN6f/jgg0 E62YfwxExz+n+HzwS8jkSuPy8e42MRUnp8+OxVuZofItKhnT9ZTns+D2Sv/zFU9In wvLzu2oi+cyp8jHfyjWQvcKOSLkuFUTL5y4wM3F48J4Ht/RrGgaHOV7igXYCAT+FF ZlbjORyXpfl231fificW3K04Ssju1SG2Crnuwy6aqNx9WhYcUpaP925LFm5OY7Ne X Bcod0AmLd+GIOSf dzlBXwkpLJyFQNwN/SBWHow+RDW7xaZKOw1A2u1O7wPNB Dj6L73AlxmyhYzeY/+lg/4YFeSbxcoJK+kGYIOqYf0YITYHpKriPyLuA/DuA3OpvO BJCOOQjl2int4Q L0r7IpTk8vv10d+mgv5KcseGSnBJS pORngTWrNIaS6uFdOOYQKz NpWWduYUyZhirMM8dqmXGVimuRXWCuCGxUZ3SFO LglaJrke= \\ \hline \end{tabular} \end{center}

\end{document}
